% Generated draft: requires mathematical and source review.
\documentclass[11pt]{article}
\usepackage[margin=25mm]{geometry}
\usepackage{fontspec}
\setmainfont{Latin Modern Roman}
\usepackage{amsmath,amssymb,mathtools,unicode-math}
\setmathfont{Latin Modern Math}
\usepackage{enumitem,needspace,xurl,hyperref}
\setlength{\parindent}{0pt}
\setlength{\parskip}{5pt}
\setlength{\emergencystretch}{3em}
\newcommand{\N}{\mathbb{N}}
\newcommand{\Z}{\mathbb{Z}}
\newcommand{\Q}{\mathbb{Q}}
\newcommand{\R}{\mathbb{R}}
\newcommand{\C}{\mathbb{C}}
\newcommand{\F}{\mathbb{F}}
\newcommand{\A}{\mathbb{A}}
\newcommand{\PP}{\mathbb{P}}
\newcommand{\E}{\mathbb{E}}
\newcommand{\eps}{\varepsilon}
\newcommand{\dd}{\,\mathrm{d}}
\begin{document}
\section*{UnsolvedMath ID: 1208}
\textbf{Source problem code:} catalogue entry
\subsection{Definitions and mathematical statement}
The phrase ``an integral curve defined by two periodic functions'' does not specify how those functions determine a curve. The supplied triage identifies the intended standard closed-curve problem as the periodic curvature-and-torsion problem, whereas the background describes coordinatewise integration. These are different mathematical questions and are retained separately.

For the classical regular space-curve interpretation, a precise setting to be confirmed is a period $T>0$ and smooth functions $\kappa,\tau:\R\to\R$ satisfying
\[
\kappa(s+T)=\kappa(s)>0,\qquad \tau(s+T)=\tau(s).
\]
Here $s$ is arc length, $\kappa$ is curvature, and $\tau$ is torsion. The Frenet curve and orthonormal frame $(\mathbf t,\mathbf n,\mathbf b)$ are determined up to an initial Euclidean motion by
\[
\gamma'=\mathbf t,\qquad
\mathbf t'=\kappa\mathbf n,\qquad
\mathbf n'=-\kappa\mathbf t+\tau\mathbf b,\qquad
\mathbf b'=-\tau\mathbf n.
\]
Normalize $\gamma(0)=0$ and $F(0)=I_3$, where $F=(\mathbf t\ \mathbf n\ \mathbf b)$ is the frame matrix. Its one-period Euclidean-motion monodromy is $(F(T),\gamma(T))$. Position closure over the selected period requires $\gamma(T)=0$; position and tangent matching additionally require $\mathbf t(T)=\mathbf t(0)$; full frame matching requires $F(T)=I_3$. Full position-and-frame closure is therefore the formal condition
\[
\gamma(T)=0,\qquad F(T)=I_3.
\]
The intended research request, according to the supplied assessment, is for explicit, usable necessary and sufficient conditions on the periodic data $\kappa,\tau$ beyond merely solving this differential system and testing its monodromy. ``Explicit'' and ``readily checked'' are not given a formal computational meaning.

This is only a specified classical setting: the stored statement does not impose regularity, positive curvature, an ambient dimension, a closure convention, or whether closure after a multiple of $T$ suffices. Curves with vanishing curvature or singularities require additional definitions, and the record does not choose such definitions. These omissions prevent a unique unconditional formulation.

For comparison, under the background's coordinatewise interpretation with continuous $T$-periodic functions $f,g$ and
\[
\eta(s)=\left(\int_0^s f(u)\,\dd u,\int_0^s g(u)\,\dd u\right),
\]
the supplied triage notes the elementary zero-mean closure criterion $\int_0^T f(u)\,\dd u=\int_0^T g(u)\,\dd u=0$. This is not the standard Efimov--Fenchel curvature-and-torsion problem.

The supplied literature review, checked 2026-08-17, labels the intended problem partially solved. It reports the formal Frenet-monodromy criterion, Bamba and Ogata's rotationally symmetric special-family solution with extensions to several singular-curve settings, and Arroyo, Garay and Mencia's explicit sufficient and nonclosure results for the planar one-curvature-function case. It reports no verified general explicit usable characterization for periodic curvature-torsion pairs. The source also mentions motivation from dynamical systems, Hamiltonian mechanics and periodic orbits. These are dated source descriptions, not a new literature assessment.
\subsection{Short English statement}
Find usable tests for when periodic curvature and torsion generate a closed space curve. The stored wording does not identify the functions as curvature and torsion or specify what kind of closure is required; coordinatewise integration gives a different, elementary zero-mean problem. The supplied assessment records monodromy and special-family progress, with the general intended characterization still open at its check date.
\subsection{Sources}
\begin{itemize}
\item UnsolvedMath local registry, permanent ID 1208, source code catalogue entry, ``The Closed Curve Problem''. Dataset attribution: Ulam AI and the UnsolvedMath Contributors; curation and metadata under CC BY 4.0.
\item K. Bamba and Y. Ogata, \emph{Closed space curves with singularities, generated by periodic curvature and torsion}, Journal of Geometry 117 (2026), article 12, \url{https://doi.org/10.1007/s00022-025-00792-3}. The supplied triage cites a rotationally symmetric special case and singular extensions.
\item Josu Arroyo, Oscar J. Garay and Jose J. Mencia, \emph{When is a periodic function the curvature of a closed plane curve?}, American Mathematical Monthly 115 (2008), 405--414, \url{https://doi.org/10.1080/00029890.2008.11920543}. The supplied triage cites planar closure and nonclosure criteria.
\item Eric W. Weisstein, \emph{Closed Curve Problem}, MathWorld, \url{https://mathworld.wolfram.com/ClosedCurveProblem.html}. The supplied assessment cites its curvature-and-torsion formulation.
\item The local literature triage records its checks as 2026-08-17; the above references were not independently visited for this draft.
\end{itemize}
\end{document}
